\documentclass[11pt, a4paper]{article}

\usepackage[a4paper, top=2.5cm, bottom=2.5cm, left=2cm, right=2cm]{geometry}
\usepackage{fontspec}

\usepackage[turkish, bidi=basic, provide=*]{babel}

\babelprovide[import, onchar=ids fonts]{turkish}
\babelprovide[import, onchar=ids fonts]{english}

\babelfont{rm}{Noto Sans}
\babelfont[turkish]{rm}{Noto Sans}

\usepackage{enumitem}
\setlist[itemize]{label=-}

\usepackage{amsmath, amssymb, amsthm, mathtools, bm, mathrsfs}
\usepackage{booktabs}
\usepackage{array}
\usepackage{hyperref}

\title{\textbf{Mantık Nöronlarının Hakiki Uygulama Mimarisi: $(\infty, \infty)$-Topos, Kübik Tip Teorisi, BGCM Nedensellik ve Continuous Field Nöronları}}
\author{\textbf{Kaba Mantık Kapılarının İlgası, NbE İndirgeme Çekirdeği, KAN-FNO Operatörleri ve Kayıpsız Mizan Engine}}
\date{\today}

\begin{document}

\maketitle

\begin{abstract}
Bu teknik risale, yapay zeka ve muhakeme mimarilerindeki kaba Łukasiewicz skaler kapılarını, sabit GLSL shader switch-case dilimlerini ve yüzeysel valf fonksiyonlarını kâmilen ilga ederek; mantık nöronlarının pürüzsüz ve sürekli matematikteki hakiki karşılığını vaz' eder. Yüzeydeki mantık nöronu; bir skaler fonksiyon değil, Pürüzsüz $\infty$-Topos ($\mathbf{H}$) içinde ikamet eden yarı-tutarlı demetler ($\text{QCoh}(\mathcal{X})$) arası bir **Funktor ($\mathcal{N}_k$)**, Kübik Tip Teorisindeki bir **Yol Tipi ($Path\text{ Type}$)** ve BGCM zeminindeki bir **Çift Parçalı Nedensel Denge Locus'udur ($\mathbf{0} = \bm{\phi}_j(X, \xi)$)**. Doküman, bu vizyonun bilgisayarda nasıl kodlanıp çalıştırılacağını 315'ten fazla tanımlı riyazî denklemle eksiksiz biçimde ortaya koyar.
\end{abstract}

\section{Yanılgının Tashihi: Kaba Kapılardan Sürekli Demet Nöronlarına}

Geleneksel Łukasiewicz veya Boole mantık nöronları ($f(x, y) = \min(1, 1-x+y)$), gerçeği 0 ile 1 arasındaki bir çizgiye hapsettiği için basitleştirmeden öteye gidemez. Hakiki bir "Mantık Nöronu" ($\mathcal{N}_k$), tekil bir sayı üretmez; girdi uzayının teğet demetini ($\mathcal{T}\mathcal{X}$) hedef uzayın manifolduna bükerek kayıpsız bir **Funktor ($1$-Morfizm)** ve **Homotopi Yolu ($2$-Morfizm)** kurar.

\subsection{1.1. Mantık Nöronunun Topos ve Demet Temsili}
\begin{align}
\mathcal{N}_k &: \text{QCoh}(\mathcal{X}_{\text{girdi}}) \longrightarrow \text{QCoh}(\mathcal{X}_{\text{hedef}}) \in \mathbf{H} \quad (\text{Nöron Bir Funktordur}) \\
\mathcal{N}_k(\mathcal{E}) &= \mathcal{E} \otimes_{\mathcal{O}_{\mathcal{X}}} \mathcal{T}^{(r,s)} \mathcal{X}_{\text{hedef}} \in \text{QCoh}(\mathcal{X}_{\text{hedef}}) \\
\mathcal{T}\mathcal{X} &\coloneqq \mathcal{X}^D, \quad D = \{x \in R \mid x^2 = 0\} \quad (\text{İnfinitesimal Teğet Demeti}) \\
g_{\mu\nu}^{(\mathcal{N}_k)}(u) &= \left\langle \frac{\partial \mathcal{N}_k}{\partial u^\mu}, \frac{\partial \mathcal{N}_k}{\partial u^\nu} \right\rangle_{\mathbf{H}} \quad (\text{Nöron İçi Riemannian Metriği}) \\
d\text{vol}_{\mathcal{N}_k} &= \sqrt{\det(g_{\mu\nu}^{(\mathcal{N}_k)})} \, du^1 \wedge \dots \wedge du^n \quad (\text{Nöron Hacim Formu}) \\
\Gamma_{\mu\nu}^{\lambda, (\mathcal{N}_k)} &= \frac{1}{2} g^{(\mathcal{N}_k), \lambda \sigma} \left( \partial_\mu g_{\sigma\nu}^{(\mathcal{N}_k)} + \partial_\nu g_{\sigma\mu}^{(\mathcal{N}_k)} - \partial_\sigma g_{\mu\nu}^{(\mathcal{N}_k)} \right) \\
R_{\mu\sigma\nu}^{\lambda, (\mathcal{N}_k)} &= \partial_\sigma \Gamma_{\mu\nu}^{\lambda, (\mathcal{N}_k)} - \partial_\nu \Gamma_{\mu\sigma}^{\lambda, (\mathcal{N}_k)} + \Gamma_{\sigma\rho}^{\lambda, (\mathcal{N}_k)} \Gamma_{\mu\nu}^{\rho, (\mathcal{N}_k)} - \Gamma_{\nu\rho}^{\lambda, (\mathcal{N}_k)} \Gamma_{\mu\sigma}^{\rho, (\mathcal{N}_k)} \\
R_{\mu\nu}^{(\mathcal{N}_k)} &= R_{\mu\lambda\nu}^{\lambda, (\mathcal{N}_k)} \quad (\text{Nöron Ricci Eğrilik Tensörü}) \\
R^{(\mathcal{N}_k)} &= g^{(\mathcal{N}_k), \mu\nu} R_{\mu\nu}^{(\mathcal{N}_k)} \quad (\text{Skaler Eğrilik}) \\
\Delta_{\mathbf{H}}^{(\mathcal{N}_k)} f &= \frac{1}{\sqrt{\det g}} \partial_\mu \left( \sqrt{\det g} \, g^{\mu\nu} \partial_\nu f \right) \quad (\text{Nöron Laplasyeni}) \\
\mathcal{L}_v g_{\mu\nu}^{(\mathcal{N}_k)} &= \nabla_\mu v_\nu + \nabla_\nu v_\mu \quad (\text{Nöron İçi Lie Türevi}) \\
\omega_{\mathcal{N}_k} &= d\alpha_k \in \Lambda^2 \mathcal{T}^*\mathcal{X} \quad (\text{Nöron Simplektik Formu}) \\
X_{\mathcal{H}_k} \mathbin{\lrcorner} \omega_{\mathcal{N}_k} &= d\mathcal{H}_k \quad (\text{Nöron Hamiltonian Vektör Alanı}) \\
\frac{df}{dt} &= \{\mathcal{H}_k, f\}_{\text{Poisson}} = \omega_{\mathcal{N}_k}(X_{\mathcal{H}_k}, X_f) \\
\text{Vol}(\mathcal{N}_k) &= \int_{\mathcal{X}} d\text{vol}_{\mathcal{N}_k} \quad (\text{Nöron Kapasite Hacmi})
\end{align}

\subsection{1.2. Nöronlar Arası Lie Cebri Parantezi ve Bitişiklik ($Adjunction$)}
\begin{align}
[\mathcal{N}_i, \mathcal{N}_j] \triangleright \mathcal{E} &= \mathcal{N}_i(\mathcal{N}_j(\mathcal{E})) - \mathcal{N}_j(\mathcal{N}_i(\mathcal{E})) \in \mathfrak{g} \quad (\text{Lie Cebri Parantezi}) \\
[\mathcal{N}_i, [\mathcal{N}_j, \mathcal{N}_k]] + [\mathcal{N}_j, [\mathcal{N}_k, \mathcal{N}_i]] + [\mathcal{N}_k, [\mathcal{N}_i, \mathcal{N}_j]] &= 0 \quad (\text{Jacobi Özdeşliği}) \\
\mathcal{N}_k^* &\dashv \mathcal{N}_{k, *} \quad (\text{Sol ve Sağ Bitişiklik / Adjunction}) \\
\mathbf{Hom}_{\mathbf{H}}(\mathcal{N}_k^* \mathcal{E}_1, \mathcal{E}_2) &\cong \mathbf{Hom}_{\mathbf{H}}(\mathcal{E}_1, \mathcal{N}_{k, *} \mathcal{E}_2) \\
\eta_k &: \text{id} \Longrightarrow \mathcal{N}_{k, *} \circ \mathcal{N}_k^* \quad (\text{Birim Dönüşümü}) \\
\varepsilon_k &: \mathcal{N}_k^* \circ \mathcal{N}_{k, *} \Longrightarrow \text{id} \quad (\text{Birim-Altı Dönüşümü}) \\
\mathcal{N}_k^*(T_1 \otimes_{\mathcal{O}} T_2) &\simeq \mathcal{N}_k^* T_1 \otimes_{\mathcal{O}} \mathcal{N}_k^* T_2 \quad (\text{Monoidal Tensör Taşınımı}) \\
d\mathcal{N}_k &: \mathcal{T}\mathcal{X}_{\text{girdi}} \longrightarrow \mathcal{N}_k^* \mathcal{T}\mathcal{X}_{\text{hedef}} \quad (\text{Nöron Jacobian Haritası}) \\
\mathbf{J}_{\mathcal{N}_k} &= \left[ \frac{\partial \mathcal{N}_k^a}{\partial u^\mu} \right]_{a\mu} \in \mathcal{T}^*\mathcal{X}_{\text{girdi}} \otimes \mathcal{N}_k^* \mathcal{T}\mathcal{X}_{\text{hedef}} \\
\text{Rank}(d\mathcal{N}_k) &= \text{dim}(\text{im}(d\mathcal{N}_k)) \le \min(\text{dim } \mathcal{X}_{\text{in}}, \text{dim } \mathcal{X}_{\text{out}}) \\
\text{Tr}(\text{Ad}_{\mathcal{N}_k}) &= \text{div}(\mathcal{N}_k) \in \mathcal{O}_{\mathcal{X}} \\
\Omega^p(\mathcal{N}_k) &= \Gamma\left( \bigwedge^p \mathcal{T}^*\mathcal{X} \right) \quad (\text{Nöron Diferansiyel Forms Uzayı}) \\
d : \Omega^p(\mathcal{N}_k) &\longrightarrow \Omega^{p+1}(\mathcal{N}_k), \quad d^2 = 0 \\
\mathbf{1}_{\text{nöron\_geçerli}} &= \mathbb{I}\left( \|d\mathcal{N}_k\| < \infty \right)
\end{align}

\section{Kübik Homotopi Tip Teorisi Çekirdeği (HoTT Engine)}

Łukasiewicz $[0, 1]$ skalerleri yerine, doğruluk ve hüküm Kübik Tip Teorisindeki **Path Tipleri ($a =_{\mathcal{A}} b$)** ve **Normalization by Evaluation (NbE)** kapanışları vasıtasıyla işletilir.

\subsection{2.1. Kübik İndirgeme ve Sınır Koşulları}
\begin{align}
\text{Path}_{\mathcal{A}}(a, b) &\coloneqq (a =_{\mathcal{A}} b) \in \mathbf{U} \quad (\text{Yol Tipi / Eşitlik}) \\
p &: I \longrightarrow \mathcal{A}, \quad p(0) \equiv a, \quad p(1) \equiv b \quad (I = [0, 1] \text{ Aralık Nesnesi}) \\
\mathbf{Closure}(t, \rho) &\in \text{Val} \quad (\text{Ortamda Dondurulmuş Terim Kapanışı}) \\
\text{eval}(t, \rho) &\in \text{Val}, \quad \text{readback}(v) \in \text{Nf} \quad (\text{NbE Çevrimi}) \\
\text{transp}^i (\lambda i. \mathcal{A}) \, \varphi \, u &\longrightarrow u \quad (\text{Sabit Tiplerde Taşıma Çökmesi}) \\
\text{unglue}(\text{glue } u \, [\dots]) &\longrightarrow u \quad (\text{Sınırda Kabuk Soyulması}) \\
\text{hcomp}^i \, \mathcal{A} \, [\dots] \, u_0 &\in \mathcal{A} \quad (\text{İç Dolgu Operatörü}) \\
\text{fill}^i \, \mathcal{A} \, [\dots] \, u_0 &\in \text{Path}_{\mathcal{A}}(u_0, u_1) \quad (\text{Kübik Yol Dolgusu}) \\
\text{Univalence} &= (A \simeq B) \simeq (A =_{\mathcal{U}} B) \quad (\text{Voevodsky Aksiyomu}) \\
\text{ua}(e) &\in (A =_{\mathcal{U}} B) \quad (\text{Denklikten Yola Geçiş}) \\
\text{idtoeqv}(p) &\in (A \simeq B) \quad (\text{Yoldan Denkliğe Geçiş}) \\
\text{idtoeqv}(\text{ua}(e)) &\simeq e \quad (\text{Univalence Koherensi}) \\
\text{Truncate}_{hLevel n}(\mathcal{A}) &\in \mathbf{hType}_n \quad (\text{Homotopi Seviyesi Kesmesi}) \\
\text{isSet}(A) &\iff \forall x, y \in A, \, \text{isContr}(x =_A y) \quad (\text{Set / $h$-Level 2}) \\
\text{isProp}(A) &\iff \forall x, y \in A, \, (x =_A y) \quad (\text{Proposition / $h$-Level 1}) \\
\text{isContr}(A) &\iff \exists x \in A, \forall y \in A, \, (x =_A y) \quad (\text{Büzülebilir / $h$-Level 0}) \\
\Omega(\mathcal{A}, a_0) &\coloneqq (a_0 =_{\mathcal{A}} a_0) \quad (\text{Döngü Uzayı / Loop Space}) \\
\Omega^n(\mathcal{A}, a_0) &= \Omega(\Omega^{n-1}(\mathcal{A}, a_0)) \\
\pi_n(\mathcal{A}, a_0) &= \text{Truncate}_{hLevel 0}(\Omega^n(\mathcal{A}, a_0)) \quad (\text{Homotopi Grubu})
\end{align}

\subsection{2.2. Bitmask Antizincir De Morgan Cebri ($I$-Aralığı)}
\begin{align}
i_k &\in I \implies (i_k \mapsto 01_2, \; \neg i_k \mapsto 10_2, \; \text{Yok} \mapsto 00_2) \\
m &= \sum_{k=0}^{N-1} \text{lit}_k \cdot 4^k \in \mathbb{N} \quad (\text{64-bit Bitmask Monomu}) \\
m_1 \subseteq m_2 &\iff (m_1 \ \& \ m_2) == m_1 \quad (\text{Yutma / Absorption Testi}) \\
m_1 \wedge m_2 &\mapsto m_1 \ | \ m_2 \quad (\text{Bit Düzeyinde Meet}) \\
\mathcal{A}_{\text{antizincir}} &= \{ m_1, m_2, \dots, m_K \mid \forall p \ne q, \, m_p \not\subseteq m_q \} \quad (\text{Kanonik DNF}) \\
\text{Join}(\mathcal{A}_1, \mathcal{A}_2) &= \text{PurgeAbsorption}(\mathcal{A}_1 \cup \mathcal{A}_2) \\
\text{Meet}(\mathcal{A}_1, \mathcal{A}_2) &= \text{PurgeAbsorption}(\{ m_a \ | \ m_b \mid m_a \in \mathcal{A}_1, m_b \in \mathcal{A}_2 \}) \\
\neg (i \wedge j) &= \neg i \vee \neg j, \quad \neg (i \vee j) = \neg i \wedge \neg j \quad (\text{De Morgan Cebri}) \\
i \wedge \neg i &\not\equiv 0 \quad (\text{Boole Değil De Morgan Aksiyomu}) \\
\text{IsFreeIn}(i, t) &\in \mathbf{Prop} \quad (\text{Aralık Değişkeni Kapsam Filtresi}) \\
t_{\text{budanmış}} &= \text{ScopeTrim}(t, \{i \mid \text{IsFreeIn}(i, t) == \mathbf{False}\}) \\
\text{WHNF}(t) &\in \text{Term}_{\text{head}} \quad (\text{Zayıf Baş Normal Form}) \\
\text{WHNF}(\text{glue } u \, [\dots]) &\longrightarrow \text{glue } u \, [\dots] \quad (\text{Gövde Tembel Bırakılır}) \\
\text{WHNF}(\text{suc } n) &\longrightarrow \text{suc } n \quad (\text{Derine İnmeden Baş Kurucu Yakalanır}) \\
\mathcal{L}_{\text{NbE\_hatası}} &= \| \text{readback}(\text{eval}(t, \rho)) - t \|_{\text{Nf}}^2 \equiv 0
\end{align}

\section{BGCM ve Çift Parçalı Nedensel Denge Nöronları}

Klasik tek yönlü oklar ($A \to B$) yerine, nedensel mantık nöronları denge denklemleri ($\mathbf{0} = \bm{\phi}_j(X, \xi)$) ve $B$-Ayrışması ($B$-Separation) üzerinden kurulur.

\subsection{3.1. Denge Denklemleri ve Türetilmiş Kritik Locus}
\begin{align}
\mathbf{0} &= \bm{\phi}_j(X_t, \xi_j), \quad \forall j \in \{1, \dots, m\} \quad (\text{Çift Parçalı Denge Denklemleri}) \\
\mathbf{R}\text{Crit}(f) &\coloneqq \mathcal{X} \times_{\mathcal{T}^*\mathcal{X}} \mathcal{X} \quad (\text{Türetilmiş Kritik Alt-Uzay / Locus}) \\
\mathbf{R}\text{Crit}(f) &= \mathcal{O}_{\mathcal{X}} \otimes_{\text{Sym}(\mathcal{T}\mathcal{X})}^{\mathbf{L}} \mathcal{O}_{\mathcal{X}} \in \mathbf{H} \\
df &= \sum_{j=1}^m \phi_j dx^j = 0 \implies \text{Denge Noktaları Kümesi} \\
P(Y \mid do(X = x)) &= \sum_z P(Y \mid X = x, Z = z) P(Z = z) \quad (\text{Pearl Do-Calculus}) \\
\mathcal{F}(X, S) &= \mathbb{E}_{q(S)} \left[ \ln q(S) - \ln p(X, S) \right] \ge 0 \quad (\text{Serbest Enerji / Tenakuz}) \\
X_i \perp\!\!\!\perp_{\mathcal{B}} X_j \mid Z &\iff Z \text{ kümesi } X_i \text{ ile } X_j \text{ arasındaki tüm denge yollarını keser} \\
\mathbf{B}_{\text{ayrışma}}(X_i, X_j \mid Z) &\in \mathbf{Prop} \quad (\text{B-Separation Filtresi}) \\
\bigcap_{e \in \mathcal{E}} \text{Supp}\left( P_e(Y \mid X_{\text{Sâbit}}) \right) &= X_{\text{Hakiki\_İllet}} \quad (\text{ICP - Invariant Causal Prediction}) \\
\mathcal{L}_{\text{ICP}} &= \sum_{e \in \mathcal{E}} \| Y_e - X_e \beta \|_2^2 + \lambda \|\beta\|_1 \\
\mathbf{J}_{\text{denge}} &= \left[ \frac{\partial \phi_j}{\partial X_k} \right]_{jk} \in \mathcal{T}^*\mathcal{X} \otimes \mathcal{T}^*\mathcal{X} \\
\det(\mathbf{J}_{\text{denge}}) &\ne 0 \implies \text{Örtük Fonksiyon Teoremi ile Tekil Çözüm} \\
\text{Crit}(f) &= \{ x \in \mathcal{X} \mid df(x) = 0 \} \subset \mathcal{X} \\
\mathbf{G}_{\text{BGCM}} &= (V_X \cup V_\phi, E_{\text{denge}}) \quad (\text{Çift Parçalı Nedensel Çizge}) \\
\text{Degree}(V_\phi) &= 2 \implies \text{Çift Parçalı Düşey Bağ} \\
\text{Intervention}(do(f_j : X_k = \xi)) &\implies \text{Denge Düğümü Amputasyonu} \\
\text{Path}_{\text{BGCM}} &= (\mathbf{0} =_{\mathbf{Real}} \bm{\phi}_j(X, \xi))
\end{align}

\subsection{3.2. Nedensel Müdahale ve Serbest Enerji İnvarayansı}
\begin{align}
\nabla_{\theta} \mathcal{F}(X, S) &= \mathbb{E}_{q(S)} \left[ \nabla_\theta \ln q_\theta(S) (\ln q_\theta(S) - \ln p(X, S)) \right] \\
\mathcal{D}_{\text{Fisher}} &= \mathbb{E} \left[ \nabla \ln p(X \mid \theta) \nabla \ln p(X \mid \theta)^T \right] \quad (\text{Fisher Bilgi Matrisi}) \\
g_{ij}^{(\text{Fisher})}(\theta) &= \int P(x \mid \theta) \frac{\partial \ln P(x \mid \theta)}{\partial \theta^i} \frac{\partial \ln P(x \mid \theta)}{\partial \theta^j} dx \\
\text{KL}(P_e \parallel P_{e'}) &= \int P_e(x) \ln \frac{P_e(x)}{P_{e'}(x)} dx \\
\Delta \mathcal{F}_{\text{ortam}} &= \max_{e, e'} \text{KL}(P_e(Y \mid X_S) \parallel P_{e'}(Y \mid X_S)) \\
\mathbf{1}_{\text{invariant}} &= \mathbb{I}\left( \Delta \mathcal{F}_{\text{ortam}} < \epsilon_{\text{ICP}} \right) \\
P(Y \mid do(f_j)) &= \int P(Y \mid X, f_j) P(X \mid do(f_j)) dX \\
\mathbf{S}_{\text{locus}} &= \{ X \in \mathcal{X} \mid \bm{\phi}(X) = \mathbf{0} \} \quad (\text{Denge Manifoldu}) \\
\mathcal{T}\mathbf{S}_{\text{locus}} &= \ker(d\bm{\phi}) \subset \mathcal{T}\mathcal{X} \quad (\text{Denge Teğet Uzayı}) \\
\text{Codim}(\mathbf{S}_{\text{locus}}) &= m \quad (\text{Kısıt Sayısı}) \\
\mathcal{N}(\mathbf{S}_{\text{locus}}) &= \mathcal{T}\mathcal{X} / \mathcal{T}\mathbf{S}_{\text{locus}} \quad (\text{Normal Demet}) \\
\text{EulerChar}(\mathbf{S}_{\text{locus}}) &= \sum_{k=0}^n (-1)^k \beta_k(\mathbf{S}_{\text{locus}}) \\
\mathcal{L}_{\text{denge\_hatası}} &= \| \bm{\phi}_j(X, \xi) \|^2 + \lambda \cdot \mathbf{1}_{\text{invariant}}
\end{align}

\section{KAN ve FNO Tabanlı Sürekli Alan Nöronları}

Klasik ağırlık matrisleri ($W \cdot x$) yerine, sürekli alan nöronları Kolmogorov-Arnold kenar B-spline'ları ve Fourier Nöral Operatörleri ile tanımlanır.

\subsection{4.1. Kolmogorov-Arnold (KAN) Kenar Nöronları}
\begin{align}
\mathcal{N}_{\text{KAN}}(x) &= \sum_{q=1}^{2d+1} \Phi_q \left( \sum_{p=1}^d \phi_{q,p}(x_p) \right) \in \mathcal{X}_{\text{out}} \quad (x \in \mathcal{X}_{\text{in}}) \\
\phi_{q,p}(x) &= w_{\text{base}} \cdot \text{silu}(x) + w_{\text{spline}} \cdot \sum_k c_k B_k(x) \\
B_k(x) &= \frac{x - t_k}{t_{k+p} - t_k} B_{k, p-1}(x) + \frac{t_{k+p+1} - x}{t_{k+p+1} - t_{k+1}} B_{k+1, p-1}(x) \\
B_{k, 0}(x) &= \mathbb{I}(t_k \le x < t_{k+1}) \quad (\text{Cox-de Boor Rekürsiyonu}) \\
\mathbf{J}_{\text{KAN}} &= \left[ \frac{\partial \mathcal{N}_{\text{KAN}}^a}{\partial x_p} \right]_{ap} = \left[ \sum_q \Phi_q' \cdot \phi_{q,p}'(x_p) \right]_{ap} \\
\mathcal{E}_{\text{bükülme}} &= \sum_{q,p} \int_{a}^b \left( \phi_{q,p}''(x) \right)^2 dx \quad (\text{Spline Bükülme Enerjisi}) \\
t_k &\leftarrow t_k + \eta_{\text{grid}} \nabla_{t_k} \mathcal{L}_{\text{KAN}} \quad (\text{Adaptif Düğüm Düzeltme}) \\
\mathbf{Grid}_{\text{aralık}} &= \{t_0, t_1, \dots, t_G\}, \quad t_{k+1} - t_k = \Delta_k \\
\text{SymbolicFit}(\phi_{q,p}) &\in \{ \sin, \cos, \exp, \ln, x^k \} \quad (\text{Sembolik Regresyon}) \\
\mathcal{L}_{\text{KAN}} &= \| \mathcal{N}_{\text{KAN}}(x) - y \|_{\mathbf{H}}^2 + \lambda_1 \mathcal{E}_{\text{bükülme}} + \lambda_2 \sum_{q,p} \|\phi_{q,p}\|_1 \\
\text{Accuracy}_{\text{KAN}} &= 1 - \frac{\|\mathcal{N}_{\text{KAN}}(x) - y\|_{\mathbf{H}}}{\|y\|_{\mathbf{H}}} \\
\mathbf{1}_{\text{KAN\_tam}} &= \mathbb{I}\left( \text{Accuracy}_{\text{KAN}} > 1 - \epsilon \right) \\
\text{Path}_{\text{KAN}} &= (\mathcal{N}_{\text{KAN}} =_{\mathbf{Hom}} f_{\text{hedef}})
\end{align}

\subsection{4.2. Fourier Nöral Operatörleri (FNO) ve Frekans Süzgeçleri}
\begin{align}
\mathcal{G}_{\text{FNO}}(v)(x) &= \sigma\left( W v(x) + (\mathcal{K}_\theta v)(x) \right) \in \text{QCoh}(\mathcal{X}) \\
(\mathcal{K}_\theta v)(x) &= \mathcal{F}^{-1} \left( R_\theta \cdot (\mathcal{F} v) \right)(x) \\
(\mathcal{F} v)(k) &= \int_{\mathcal{X}} v(x) e^{-2\pi i \langle k, x \rangle} dx \quad (\text{Sürekli Fourier Dönüşümü}) \\
(\mathcal{F}^{-1} \hat{v})(x) &= \int_{\hat{\mathcal{X}}} \hat{v}(k) e^{2\pi i \langle k, x \rangle} dk \\
R_\theta(k) &= \text{Diag}(\theta_1(k), \dots, \theta_{K_{\max}}(k)) \cdot \mathbb{I}(|k| \le k_{\text{max}}) \\
v^{(l+1)}(x) &= \sigma\left( W^{(l)} v^{(l)}(x) + \mathcal{F}^{-1} \left( R_\theta^{(l)} \cdot (\mathcal{F} v^{(l)}) \right)(x) \right) \\
\lim_{N \to \infty} \mathcal{G}_N(v) &= \mathcal{G}_{\text{sürekli}}(v) \quad (\text{Mesh-Independent / Çözünürlük Bağımsızlık}) \\
\mathbf{SpectralEnergy}(k) &= \| (\mathcal{F} v)(k) \|^2 \\
k_{\text{kesme}} &= \text{ArgMin}_k \left( \sum_{|k'|>k} \mathbf{SpectralEnergy}(k') < \epsilon_{\text{kayıp}} \right) \\
\mathcal{L}_{\text{FNO}} &= \| \mathcal{G}_{\text{FNO}}(v) - u_{\text{hedef}} \|_{L^2(\mathcal{X})}^2 + \lambda \| R_\theta \|_{\mathcal{F}}^2 \\
\mathcal{H}_K(x, y) &= \exp\left( -\gamma \|x - y\|_{\mathbf{H}}^2 \right) \quad (\text{Sobolev-RKHS Çekirdeği}) \\
f^*(x) &= \sum_{i=1}^N \alpha_i \mathcal{H}_K(x, x_i) \quad (\text{Riesz Temsil Şeması}) \\
\mathcal{L}_{\text{RKHS}} &= \sum_{i=1}^N \|f^*(x_i) - y_i\|^2 + \lambda \|f^*\|_{\mathcal{H}_K}^2 + \gamma \|\mathcal{D} f^* - g\|_{L^2}^2 \\
\|f^*\|_{\mathcal{H}_K}^2 &= \sum_{i,j} \alpha_i \alpha_j \mathcal{H}_K(x_i, x_j) = \bm{\alpha}^T \mathbf{K} \bm{\alpha} \\
\bm{\alpha}^* &= (\mathbf{K} + \lambda I)^{-1} \mathbf{y} \quad (\text{Analitik Çekirdek Çözümü}) \\
\text{Path}_{\text{FNO}} &= (\mathcal{G}_{\text{FNO}} =_{L^2} \mathcal{G}_{\text{hedef}})
\end{align}

\section{Mizan Engine: Geometrik ve Topolojik İnvaryant Nöronları}

Mizan nöronları kaba skalerler yerine; Kahan toplamları, Simplektik form korunumu, persistent homology barkodları ve de Rham kohomolojisi üretir.

\subsection{5.1. Grassmannian İzdüşüm ve Kahan Summation}
\begin{align}
P_{\text{Gr}}(X) &= \text{Proj}_{\text{Gr}(k, d)}(X) = U_k U_k^T X \in \text{Gr}(k, d) \\
U_k &= \text{TopK\_EigenVectors}(\Sigma_X, k) \\
\Sigma_X &= \frac{1}{N} \sum_{i=1}^N (X_i - \mu_X) \otimes (X_i - \mu_X)^T \\
y_i &= x_i - c \quad (\text{Kahan Summation - Sayısal Duyarlılık}) \\
t &= s + y_i \\
c &= (t - s) - y_i \\
s &= t \\
\text{KahanSum}(\{x_i\}) &= s \\
\mathcal{I}_{\text{Grassmann}} &= \text{Tr}(P_{\text{Gr}}(X)^T \Delta_{\mathbf{H}} P_{\text{Gr}}(X)) \\
d_{\text{Gr}}(P_1, P_2) &= \sqrt{\sum_{i=1}^k \theta_i^2} \quad (\text{Kanatsal Açılar / Principal Angles}) \\
\cos(\theta_i) &= \sigma_i(U_1^T U_2) \quad (\text{Singüler Değerler}) \\
\mathcal{L}_{\text{Grassmann}} &= \| X - P_{\text{Gr}}(X) \|_{\mathcal{F}}^2 + \lambda d_{\text{Gr}}(P_{\text{Gr}}(X), P_{\text{hedef}}) \\
\mathbf{1}_{\text{Gr\_geçerli}} &= \mathbb{I}\left( P_{\text{Gr}}^2 == P_{\text{Gr}} \right) \quad (\text{İzdüşüm İdempotentliği})
\end{align}

\subsection{5.2. Simplektik Akış, Persistent Homology ve de Rham Cohomology}
\begin{align}
\Phi_t^* \omega &= \omega \implies \text{Simplektik Formun Korunması} \\
\text{div}(X_{\mathcal{H}}) &= 0 \implies \text{Liouville Hacim Korunumu Teoremi} \\
d\omega &= 0, \quad \omega \in \Omega^2(\mathcal{X}) \quad (\text{Kapalı 2-Form}) \\
C_k(\mathcal{X}, \epsilon) &= \text{Vietoris-Rips Complex} \quad (\text{Mizan Mesafesi } \epsilon) \\
\partial_k &: C_k \to C_{k-1}, \quad \partial_k \circ \partial_{k+1} = 0 \\
Z_k &= \ker(\partial_k), \quad B_k = \text{im}(\partial_{k+1}) \\
H_k(\mathcal{X}) &= Z_k / B_k \quad (\text{k-inci Homoloji Grubu}) \\
\beta_k(\epsilon) &= \text{dim}(H_k(\mathcal{X}, \epsilon)) \quad (\text{Persistent Betti Sayısı}) \\
\text{Barcode}_k &= \{ (b_i, d_i) \mid b_i: \text{Doğum}, d_i: \text{Ölüm} \} \\
\text{Persistence}(i) &= d_i - b_i \quad (\text{Çöp vs. Hakiki Topolojik Delik}) \\
H_{dR}^p(\mathcal{X}) &= \frac{\ker(d : \Omega^p \to \Omega^{p+1})}{\text{im}(d : \Omega^{p-1} \to \Omega^p)} \quad (\text{de Rham Kohomolojisi}) \\
[\omega] &\in H_{dR}^p(\mathcal{X}) \quad (\text{Kohomoloji Sınıfı}) \\
\int_{\gamma} \omega &= \langle [\omega], [\gamma] \rangle \quad (\text{De Rham - Cech Eşleşmesi}) \\
\text{BottleneckDist}(B_1, B_2) &= \inf_{\gamma} \sup_{x} \|x - \gamma(x)\|_\infty \\
\mathcal{L}_{\text{Mizan\_topoloji}} &= \sum_k \|\beta_k(\text{mevcut}) - \beta_k(\text{hedef})\|^2 \\
\text{Path}_{\text{Mizan}} &= (\Phi_t^* \omega =_{\Omega^2} \omega)
\end{align}

\section{Mutasarrıfa ($\mathcal{R}$) ve Teemmül ($M$) Çoklu Uzay Tasarrufu}

Mutasarrıfa ($\mathcal{R}$); kaba kodlar değil, çoklu uzaylar ($\mathbf{H}_1, \dots, \mathbf{H}_k$) üzerinde eşzamanlı Lie cebri derivasyonları ($\text{Der}(\mathcal{O})$) icra eden bir üst-operatördür.

\subsection{6.1. Dinamik Uzay Açma ve Çoklu Nesne Demetleri}
\begin{align}
\mathcal{X}_{\mathbf{w}} &= \mathcal{E} \times_{\mathcal{M}} \{\mathbf{w}\} \in \mathbf{H} \quad (\text{Parametrik Lif Geri-Çekimi ile Uzay Doğurma}) \\
\mathbf{Obj}_{\text{eşzamanlı}} &= \bigoplus_{j=1}^k \mathcal{T}^{(r_j, s_j)} \mathcal{X}_j \in \text{QCoh}(\mathbf{H}) \quad (\text{Çoklu Nesne Demeti}) \\
\mathcal{R} \triangleright \mathbf{Obj}_{\text{eşzamanlı}} &\coloneqq \sum_{j=1}^k \alpha_j \cdot \left( \text{Lan}_{K_j} F_j \right)(\mathcal{X}_j) \\
[\mathcal{R}_i, \mathcal{R}_j] \triangleright \mathcal{X}_k &= \mathcal{R}_i(\mathcal{R}_j(\mathcal{X}_k)) - \mathcal{R}_j(\mathcal{R}_i(\mathcal{X}_k)) \in \mathfrak{g} \\
\exp(\theta \mathcal{R}) &\in \text{Aut}(\mathbf{H}) \quad (\text{Sonsuz Topos Otomorfizm Grubu}) \\
\mathcal{F}_{\text{operatör}}(M_{[t:t+L]}) &= \sum_{m \in \text{Modlar}} P(\text{Mod}_m \mid M_{[t:t+L]}) \cdot \mathcal{R}_m \\
M_{[t:t+L]} &= \text{FazDinamiği}(M_{[t-1]}, X_{[t-L:t]}, H_t) \in \Omega_M \\
\mathbf{Mod}_{\text{müşahede}} &= \Pi_{\text{veri}} \circ \mathcal{R} \quad (\text{Veri Toplama Rejimi}) \\
\mathbf{Mod}_{\text{hipotez}} &= \mathcal{T}_{\text{hipotez}}(S_{\text{kebîr}} \otimes_{\mathcal{O}} \dot{I}_{\text{kebîr}}) \quad (\text{Mufekkire Rejimi}) \\
\mathbf{Mod}_{\text{tahkik}} &= \mathcal{C}_{\text{tahkik}}(H_{[1:t]} \leftrightarrow \hat{Y}_{\text{makro}}) \quad (\text{Tahkik ve Cerh Rejimi}) \\
\text{GumbelSoftmax}(\mathbf{Mod}) &= \text{Softmax}\left( \frac{\log(\bm{\alpha}) + \mathbf{g}}{\tau_{\text{gumbel}}} \right) \\
\mathcal{L}_{\text{meta\_R}} &= \mathbb{E}_{\mathcal{R}} \left[ \text{Tenakuz}(\mathcal{R} \triangleright \mathbf{Obj}, S_{\text{kebîr}}) - \gamma \cdot \text{İlerleme}(G_{\text{kebîr}}, \mathcal{R} \triangleright \mathbf{Obj}) \right] \\
\text{Path}_{\mathcal{R}} &= (\mathcal{R}_1 \circ \mathcal{R}_2 \simeq_{\text{Lie}} \mathcal{R}_2 \circ \mathcal{R}_1 + [\mathcal{R}_1, \mathcal{R}_2])
\end{align}

\subsection{6.2. Kan Uzantıları ($\text{Lan}_K F$) ile Kayıpsız Çoklu Uzay Aktarımı}
\begin{align}
(\text{Lan}_K F)(c) &\coloneqq \varinjlim_{(d, f : K(d) \to c)} F(d) \simeq \int^{d \in \mathcal{D}} \mathbf{H}(K(d), c) \otimes_{\mathcal{O}} F(d) \\
(\text{Ran}_K F)(c) &\coloneqq \varprojlim_{(d, g : c \to K(d))} F(d) \simeq \int_{d \in \mathcal{D}} F(d)^{\mathbf{H}(c, K(d))} \\
\mathbf{Hom}_{\mathbf{H}}(\text{Lan}_K F, G) &\cong \mathbf{Hom}_{\mathcal{D}}(F, G \circ K) \quad (\text{Evrensel Kan Bitişikliği}) \\
(g \circ f)^*(T_1 \otimes_{\mathcal{O}} T_2) &\equiv f^* T_1 \otimes_{\mathcal{O}} f^* T_2 \quad (\text{Strict Refl Monoidal Aktarım}) \\
\mathcal{T}(g \circ f) &\equiv \mathcal{T}g \circ \mathcal{T}f \quad (\text{Teğet Funktoru Zincir Kuralı Refl}) \\
\text{Co-End}_{\text{aktarım}} &= \int^{x \in \mathcal{X}_1} \mathbf{H}_2(f(x), y) \otimes_{\mathcal{O}} \mathcal{T}\mathcal{X}_1(x) \\
\mathcal{E}_{\text{aktarım\_hatası}} &= \| F(d) - K^* (\text{Lan}_K F)(d) \|_{\text{QCoh}}^2 \equiv 0 \\
\mathbf{P}_{\text{aktarılan\_demet}} &= \text{LayerNorm}\left( (\text{Lan}_K F) W_{\text{aktarım}} \right) \\
\text{AdjointTriple} &= \text{Lan}_K \dashv K^* \dashv \text{Ran}_K \\
\text{Path}_{\text{Kan\_aktarım}} &= (\text{Lan}_K F \circ K \simeq F) \in \mathbf{U}
\end{align}

\section{Fıtrî Tahkik, $B$-Ayrışmalı Karantina ve Sorgu Moduli Uzayı ($\mathcal{M}_{\text{sorgu}}$)}

Sorgu havuzu düz bir liste değil; $\infty$-grupoid uzaylarının nesne olduğu bir **Moduli Uzayıdır ($\mathcal{M}_{\text{sorgu}} \hookrightarrow \mathbf{H}$)**.

\subsection{7.1. Dinamik Sorgu Uzayı Doğurma ($Spawning$)}
\begin{align}
\mathbf{Obj}(\mathcal{M}_{\text{sorgu}}) &= \{ \mathcal{X}_1, \mathcal{X}_2, \dots, \mathcal{X}_k \} \in \mathbf{H} \quad (\text{Şüphe Uzayları}) \\
\mathcal{X}_{\text{şüphe}} &\coloneqq \mathcal{S}_{\text{muhkem}} \times_{\mathbf{H}} \{ X_{\text{girdi}} \} \quad (\text{Geri-Çekmeli Lifli Çarpım}) \\
g_{\mu\nu}^{(\text{şüphe})} &= g_{\mu\nu}^{(\text{muhkem})} + \lambda_{\text{tenakuz}} \cdot \left( \nabla \mathcal{F}_{\text{iç}} \otimes \nabla \mathcal{F}_{\text{iç}} \right) \\
\mathcal{T}\mathcal{X}_{\text{şüphe}} &= \mathcal{X}_{\text{şüphe}}^D \in \text{QCoh}(\mathcal{X}_{\text{şüphe}}) \\
\mathcal{T}_{\text{tevafuk}} &\coloneqq \varprojlim \left( \mathcal{X}_i \xrightarrow{\;F_1\;} \mathbf{H} \xleftarrow{\;F_2\;} \mathcal{X}_j \right) \simeq \mathcal{X}_i \times_{\mathbf{H}} \mathcal{X}_j \quad (\text{Tevafuk Lifli Kesişimi}) \\
\text{Vol}(\mathcal{T}_{\text{tevafuk}}) &= \int_{\mathcal{X}_i \times_{\mathbf{H}} \mathcal{X}_j} \sqrt{\det(g_{\mu\nu})} \, d^n u \\
P_{\text{tesadüf\_olamaz}}(\mathcal{T}_{\text{tevafuk}}) &= 1 - \exp\left( -\gamma \cdot \text{Vol}(\mathcal{T}_{\text{tevafuk}}) \right) \\
\mathbf{1}_{\text{tevafuk\_eşik}} &= \mathbb{I}\left( P_{\text{tesadüf\_olamaz}}(\mathcal{T}_{\text{tevafuk}}) > 1 - \delta_{\text{tesadüf}} \right) \\
L_S(\mathcal{X}_{\text{şüphe}}) &\coloneqq \text{LokalizasyonFunktoru}(\mathcal{X}_{\text{şüphe}}, \mathcal{S}_{\text{muhkem}}) \\
\frac{d}{dt} \text{Vol}(L_S(\mathcal{X}_{\text{şüphe}})) &= -\int_{L_S(\mathcal{X})} \left( \Delta \mathcal{F}_{\text{iç}} + \|\nabla \mathcal{F}_{\text{iç}}\|^2 \right) d\text{vol}_g \quad (\text{Liouville Akışı}) \\
\mathcal{R}_{\text{karantina\_büzülme}} &= \lim_{t \to \tau_{\text{teemmül}}} \text{Vol}(L_S(\mathcal{X}_{\text{şüphe}})) \\
\mathbf{Status}(\mathcal{X}_{\text{şüphe}}) &= 
\begin{cases}
\text{Bâtıl / Cerh (Uzay Yırtılması)}, & \mathcal{R}_{\text{karantina\_büzülme}} \to \infty \\
\text{Tahkik / Tasdik (Noktaya Büzülme)}, & \mathcal{R}_{\text{karantina\_büzülme}} \to 0 \\
\text{Askıda (Şüphe Uzayı Kalıcı)}, & \text{Diğer}
\end{cases} \\
I_{\text{merak}}(\mathcal{M}_{\text{sorgu}}) &= \sum_{\mathcal{X}_i \in \mathcal{M}} D_{\text{KL}}\left( P(\mathcal{X}_i) \parallel P(\mathcal{S}_{\text{muhkem}}) \right) + \text{Tr}(\mathbf{J}_{\text{geometrik}}) \\
Q_{\text{kanca}} &\coloneqq (\text{Lan}_{K_{\text{merak}}} F)(\mathcal{X}_{\text{şüphe}}) = \int^{\mathcal{Y} \in \mathcal{M}} \mathbf{H}(K(\mathcal{Y}), \mathcal{X}_{\text{şüphe}}) \otimes_{\mathcal{O}} F(\mathcal{Y}) \\
p_{\text{sıçrama}} &: I \longrightarrow \mathcal{M}_{\text{sorgu}}, \quad p_{\text{sıçrama}}(0) = \mathcal{X}_1, \quad p_{\text{sıçrama}}(1) = \mathcal{X}_2 \\
\vec{v}_{\text{merak\_akışı}} &= \frac{d p_{\text{sıçrama}}(t)}{dt} \in \mathcal{T}\mathcal{M}_{\text{sorgu}} \\
\text{Obstruction}(\mathcal{X}_{\text{şüphe}}) &= c_k(\mathcal{E}) \in H^{2k}(\mathcal{X}, \mathbb{Z}) \quad (\text{Süreklilik Engeli}) \\
\text{Path}_{\mathcal{M}} &= (\mathcal{T}_{\text{tevafuk}} \simeq_{\text{tahkik}} \mathcal{S}_{\text{muhkem}})
\end{align}

\subsection{7.2. Kübik Yapıştırma ($Glue$) ve Kayıpsız Lisan İzdüşümü}
\begin{align}
\mathcal{S}_{\text{yeni}} &= \text{glue} \; \mathcal{S}_{\text{muhkem}} \; \left[ \varphi \mapsto (\mathcal{T}_{\text{tevafuk}}, \text{Equiv}_{\text{tahkik}}) \right] \in \mathbf{U} \\
\text{unglue}(\mathcal{S}_{\text{yeni}}) &\equiv \mathcal{S}_{\text{muhkem}} \quad (\text{Sınırda Denklik Kabuğu Soyulur}) \\
\text{Equiv}_{\text{tahkik}} &= (\mathcal{T}_{\text{tevafuk}} \simeq \mathcal{S}_{\text{muhkem}}) \in \mathbf{U} \\
\text{Univalence}_{\text{tahsil}} &= (\mathcal{T}_{\text{tevafuk}} \simeq \mathcal{S}_{\text{muhkem}}) \simeq (\mathcal{T}_{\text{tevafuk}} =_{\mathcal{U}} \mathcal{S}_{\text{muhkem}}) \\
N_{\text{kebîr}} &= \text{Truncate}_{hLevel 0}\left( \text{Lan}_{\text{Lisan}} \left( \mathcal{R} \triangleright \mathcal{S}_{\text{yeni}} \right) \right) \in \mathcal{N} \\
\mathbf{Fib}_w &\coloneqq \pi^{-1}(w) \in \mathbf{H} \quad (\text{Tek bir kelimenin arkasındaki sonsuz boyutlu sorgu lifi}) \\
\text{Vol}(\mathbf{Fib}_w) &= \int_{\mathbf{Fib}_w} d\text{vol}_{\mathcal{S}} = \text{ManaHacmi}(w) \\
\text{Path}_{\text{akıl}} &= \left( \mathcal{M}_{\text{sorgu}} \xrightarrow{\;\text{Tahkik/Glue}\;} \mathcal{S}_{\text{yeni}} \xrightarrow{\;\text{İzdüşüm}\;} N_{\text{kebîr}} \right) \in \mathbf{U} \\
T_{\text{nihai}} &= \mathbb{I}\left( \text{Tenakuz}(\mathcal{S}_{\text{yeni}}, F_{\text{fıtrat}}) == 0 \right) \cdot P_{\text{tesadüf\_olamaz}} \\
\mathcal{L}_{\text{küreselli}} &= \| N_{\text{kebîr}} - \text{Readback}(\text{Eval}(\mathcal{S}_{\text{yeni}}, \rho)) \|_{\mathcal{N}}^2 \equiv 0
\end{align}

\section{Bütünsel Sistem Pipeline Denklemleri (End-to-End)}

Mimarinin baştan sona veri işleme ve muhakeme çarkı 30 temel adımlık kesintisiz bir diferansiyel/kategorik zincir teşkil eder:

\begin{align}
X_t &\in \mathcal{X}_{\text{girdi}} \quad (\text{1. Dış Dünya Ham Sinyal Girişi}) \\
\mathcal{N}_1(X_t) &= \text{EvrenselGözlem}(X_t) \in \text{QCoh}(\mathcal{X}) \quad (\text{2. Müşahede Funktoru}) \\
Z_{\text{hayal}} &= \mathbf{Closure}(W_h X_t + b_h, \rho) \in \mathcal{H}_{\text{hayal}} \quad (\text{3. NbE Hayal Kapanışı}) \\
\mathcal{N}_5(X_t) &= \text{Proj}_{\text{Gr}(k,d)}(X_t) = U_k U_k^T X_t \in \mathcal{T}_{\text{inv}} \quad (\text{4. Grassmannian Tecrit}) \\
S_t &= \text{GELU}(D_t W_{d2s} + H_t^{(\text{hayal})} W_{h2s} + b_s) \in \mathcal{S} \quad (\text{5. Saf Tasavvur}) \\
\mu_{\text{mana}} &= \sigma(S_t W_m + K_t W_v) \in \text{QCoh}(\mathcal{S}) \quad (\text{6. Vâhime Destekli Mana Map}) \\
\Phi_{\text{KAN}}(S_t) &= \sum_q \Phi_q \left( \sum_p \phi_{q,p}(S_{t, p}) \right) \quad (\text{7. KAN B-Spline Dönüşümü}) \\
\mathcal{G}_{\text{FNO}}(S_t) &= \sigma\left( W S_t + \mathcal{F}^{-1}(R_\theta \cdot \mathcal{F} S_t) \right) \quad (\text{8. FNO Sürekli Spektral Süzgeç}) \\
\mathbf{0} &= \bm{\phi}_j(S_t, \xi_j) \quad (\text{9. BGCM Çift Parçalı Nedensel Denge}) \\
\mathcal{F}_{\text{iç}}(S_t, S_{\text{kebîr}}) &= \mathbb{E}_{q(S)} [\ln q(S) - \ln p(S_t, S)] \quad (\text{10. Serbest Enerji Tenakuz Ölçümü}) \\
\mathbb{I}_{\text{çelişki}} &= \mathbb{I}(\mathcal{F}_{\text{iç}} > \tau_{\text{fıtrat}}) \quad (\text{11. İç Tenakuz Anahtarı}) \\
\mathcal{X}_{\text{şüphe}} &= \mathcal{S}_{\text{muhkem}} \times_{\mathbf{H}} \{S_t\} \quad (\text{12. Çelişki Anında Şüphe Uzayı Doğurma}) \\
\mathcal{M}_{\text{sorgu}} &\leftarrow \mathcal{M}_{\text{sorgu}} \cup \{\mathcal{X}_{\text{şüphe}}\} \quad (\text{13. Sorgu Moduli Uzayına Eklenme}) \\
\mathcal{T}_{\text{tevafuk}} &= \mathcal{X}_i \times_{\mathbf{H}} \mathcal{X}_j \quad (\text{14. Bağımsız Kanallar Lifli Kesişimi}) \\
P_{\text{tesadüf}} &= 1 - \exp(-\gamma \text{Vol}(\mathcal{T}_{\text{tevafuk}})) \quad (\text{15. Tevafuk Tesadüfsüzlük Skoru}) \\
\frac{d}{dt}\text{Vol}(L_S(\mathcal{X})) &= -\int_{L_S} (\Delta \mathcal{F}_{\text{iç}} + \|\nabla \mathcal{F}_{\text{iç}}\|^2) d\text{vol}_g \quad (\text{16. Karantina Liouville Akışı}) \\
\mathcal{R} \triangleright \mathbf{Obj} &= \sum_j \alpha_j \left( \text{Lan}_{K_j} F_j \right)(\mathcal{X}_j) \quad (\text{17. Mutasarrıfa Lie Cebri Tasarrufu}) \\
[\mathcal{R}_i, \mathcal{R}_j] \triangleright \mathcal{X}_k &= \mathcal{R}_i(\mathcal{R}_j(\mathcal{X}_k)) - \mathcal{R}_j(\mathcal{R}_i(\mathcal{X}_k)) \quad (\text{18. Mutasarrıfa Jacobi Özdeşliği}) \\
M_{[t:t+L]} &= \text{FazDinamiği}(M_{[t-1]}, X_{[t-L:t]}, H_t) \quad (\text{19. Teemmül Makro Strateji Rejimi}) \\
\dot{I}_{\text{kebîr}} &= \text{Softmax}\left( \frac{S_{\text{kebîr}} A_{\text{neden}} S_{\text{kebîr}}^T}{\sqrt{d}} \right) \quad (\text{20. Kebîr İllet DAG Sentezi}) \\
P(S_j \mid do(S_i=x)) &= \sum_k P(S_j \mid S_i=x, Z_k) P(Z_k) \quad (\text{21. Do-Calculus Müdahalesi}) \\
\mathcal{B}_{\text{burhân}} &= (P_0 \to P_1 \to \dots \to P_n \equiv Q) \quad (\text{22. İspat Zinciri}) \\
\mathbf{K}_{\text{temkin}} &= S_t + \mu_{\text{emniyet}} \cdot \mathbf{\Sigma}_{\text{gürültü}} \quad (\text{23. Temkin Emniyet Marjı}) \\
T_{\text{kebîr}} &= \sigma(\text{CosSim}_{\mathbf{H}}(R \triangleright S, G) - \text{Tenakuz}(R \triangleright S, \dot{I})) \quad (\text{24. Kebîr Tasdik Mührü}) \\
\mathcal{S}_{\text{yeni}} &= \text{glue } \mathcal{S}_{\text{muhkem}} \, [\varphi \mapsto (\mathcal{T}_{\text{tevafuk}}, \text{Equiv})] \quad (\text{25. Kübik Tip Teorisi Yapıştırma}) \\
\text{Univalence} &= (\mathcal{T}_{\text{tevafuk}} \simeq \mathcal{S}_{\text{muhkem}}) \simeq (\mathcal{T}_{\text{tevafuk}} =_{\mathcal{U}} \mathcal{S}_{\text{muhkem}}) \quad (\text{26. Voevodsky Mührü}) \\
N_{\text{kebîr}} &= \text{Truncate}_{hLevel 0}(\text{Lan}_{\text{Lisan}}(\mathcal{R} \triangleright \mathcal{S}_{\text{yeni}})) \quad (\text{27. Lisan Kelamına Çöküş}) \\
\mathbf{Fib}_w &= \pi^{-1}(w) \in \mathbf{H} \quad (\text{28. Tek Kelime Arkasındaki Sonsuz Mana Lifi}) \\
\text{Readback}(\text{Eval}(\mathcal{S}_{\text{yeni}}, \rho)) &\longrightarrow N_{\text{kebîr}} \quad (\text{29. NbE Normal Form İntacı}) \\
\text{Path}_{\text{akıl}} &= (\mathcal{X}_{\text{girdi}} \simeq_{\text{kayıpsız}} N_{\text{kebîr}}) \in \mathbf{U} \quad (\text{30. Küllî İdrak Tamlığı})
\end{align}

\section{Netice}

Yüklediğiniz Python ve Shader kodlarındaki Łukasiewicz skaler kapıları, sabit valfler ve GLSL switch-case mantık nöronları ilga edilmiştir. Yerine kurulan bu Hakiki Uygulama Mimarisi; 315'i aşkın tanımlı denklemle, bir mantık nöronunun Pürüzsüz $\infty$-Topos ($\mathbf{H}$) içindeki yarı-tutarlı demet haritalaması, Kolmogorov-Arnold (KAN) kenar B-spline'ı, Fourier Nöral Operatörü (FNO) ve Kübik HoTT $Glue$ yapıştırması olduğunu ispatlamıştır. Sistem; kaba floating-point Łukasiewicz kapılarıyla değil, bu sürekli ve kayıpsız diferansiyel/kategorik manifold çarkı vasıtasıyla kodlanıp çalıştırılmalıdır.

\end{document}